\topic{Geometri Vektor Eigen}
\index{geometri vektor eigen|(}
--Lihat Topik tentang Geometri Transformasi Linear---

Karakterisasi transformasi linear dalam operasi-operasi elementer berguna dalam
beberapa hal (misalnya, kita dapat dengan mudah melihat bahwa garis dipetakan ke
garis karena proyeksi, dilatasi, refleksi, dan geseran masing-masing memetakan
garis ke garis). Namun, ketika suatu pemetaan dinyatakan sebagai komposisi dari
banyak operasi kecil---betapapun sederhananya---deskripsi itu kurang ideal.
Kita mengakhiri pembahasan dengan cara lain yang lebih menyeluruh untuk
menggambarkan pengaruh geometris transformasi pada $\Re^2$.

Gambar-gambar dalam pembahasan itu hanya memperlihatkan aksi pemetaan pada satu
atau dua anggota domain. Kita mengetahui bahwa suatu transformasi dideskripsikan
sepenuhnya oleh aksinya pada suatu basis. Jadi, secara ketat, untuk
mendeskripsikan transformasi pada $\Re^2$ cukup dijelaskan ke mana kedua vektor
dari sebarang basis dipetakan. Meskipun demikian, gambar-gambar itu tampaknya
belum menyampaikan banyak intuisi geometris. Dapatkah kita memperjelas geometri
suatu pemetaan linear dengan menambahkan informasi, tetapi tidak sedemikian
banyak hingga gambarnya menjadi membingungkan?

Suatu transformasi pada $\Re^2$ mengirimkan garis yang melalui titik asal ke
garis yang melalui titik asal. Jadi, dua titik pada garis $y=k_1x$ sama-sama
akan dikirimkan ke suatu garis, misalnya $y=k_2x$. Tinjau dua titik demikian.
Salah satunya merupakan kelipatan yang lain, sehingga titik kedua dapat ditulis
sebagai $r$ kali titik pertama untuk suatu skalar $r$.
\begin{equation*}
  \colvec{x \\ k_1x}\quad\mbox{dan}\quad\colvec{(rx) \\ k_1(rx)}
\end{equation*}
Bandingkan bayangan keduanya.
\begin{equation*}
  \begin{pmatrix}
    a  &c  \\
    b  &d
  \end{pmatrix}
  \colvec{x \\ k_1x}
  =
  \colvec{ax+ck_1x \\ bx+dk_1x}
  \qquad
  \begin{pmatrix}
    a  &c  \\
    b  &d
  \end{pmatrix}
  \colvec{(rx) \\ k_1(rx)}
  =
  \colvec{a(rx)+ck_1(rx) \\ b(rx)+dk_1(rx)}
\end{equation*}
Vektor kedua adalah $r$ kali vektor pertama, dan bayangan vektor kedua adalah
$r$ kali bayangan vektor pertama. Transformasi tersebut tidak hanya
mempertahankan kolinearitas kedua vektor, tetapi juga mempertahankan skala
relatif keduanya. Dengan kata lain, suatu transformasi memperlakukan titik-titik
pada garis yang melalui titik asal secara seragam. Untuk mendeskripsikan pengaruh
pemetaan pada seluruh garis, kita hanya perlu mendeskripsikan pengaruhnya pada
satu titik tak nol di garis itu.

Karena setiap titik dalam ruang terletak pada suatu garis yang melalui titik
asal, untuk memahami aksi transformasi linear pada $\Re^2$ cukup dipilih satu
titik dari setiap garis yang melalui titik asal (misalnya titik yang terletak
pada setengah lingkaran satuan bagian atas), lalu diperlihatkan pengaruh pemetaan
pada himpunan titik tersebut.

Berikut gambar demikian untuk suatu dilatasi sederhana.
\begin{center}
  \psset{xunit=8pt,yunit=8pt,linewidth=.4pt,arrowsize=.5pt 2.5,arrowinset=.1,
         arrowlength=1.75}
  \pspicture(-3,-1)(3,2)
    \psaxes[labels=none,linecolor=weakgray,tickstyle=bottom,
        ticks=all,ticksize=1.5pt]{<->}(0,0)(-2.75,-.75)(2.75,1.75)
     \parametricplot[plotstyle=curve,linewidth=.6pt]%
       {0}{180}{t cos t sin}
  \endpspicture
  \qquad\raisebox{12pt}{\begin{tabular}[t]{c}
                         $\longrightarrow$ \\
                         {\scriptsize
                          $\colvec{x \\ y} \mapsto \colvec{2x \\ y}$}
                       \end{tabular}}\qquad
  \pspicture(-3,-1)(3,2)
    \psaxes[labels=none,linecolor=weakgray,tickstyle=bottom,
        ticks=all,ticksize=1.5pt]{<->}(0,0)(-2.75,-.75)(2.75,1.75)
     \parametricplot[plotstyle=curve,linewidth=.6pt]%
       {0}{180}{2 t cos mul t sin}
  \endpspicture
\end{center}
Di bawah ini, pemetaan yang sama ditampilkan dengan lingkaran dan bayangannya
yang ditumpangtindihkan.
\begin{center}
  \psset{xunit=12pt,yunit=12pt,linewidth=.4pt,arrowsize=.5pt 2.5,arrowinset=.1,
         arrowlength=1.75}
  \pspicture(-3,-1)(3,2)
    \psaxes[labels=none,linecolor=weakgray,tickstyle=bottom,
        ticks=all,ticksize=1.5pt]{<->}(0,0)(-2.75,-.75)(2.75,1.75)
     \parametricplot[plotstyle=curve,linecolor=preimagegray,linewidth=.6pt]%
       {0}{180}{t cos t sin}
        \psline[linecolor=preimagegray]{->}(-1,0)(-2,0)
        \psline[linecolor=preimagegray]{->}(-.707,.707)(-1.414,.707)
        \psline[linecolor=preimagegray]{->}(.707,.707)(1.414,.707)
        \psline[linecolor=preimagegray]{->}(1,0)(2,0)
     \parametricplot[plotstyle=curve,linewidth=.6pt]%
       {0}{180}{2 t cos mul t sin}
  \endpspicture
\end{center}
Geometrinya kini tentu lebih jelas. Sebagai contoh, kita dapat melihat bahwa
beberapa garis yang melalui titik asal dipetakan ke dirinya sendiri:~sumbu-$x$
dipetakan ke sumbu-$x$, dan sumbu-$y$ dipetakan ke sumbu-$y$.

Berikut adalah pembalikan yang ditampilkan sebelumnya, kini dengan lingkaran
dan bayangannya ditumpangtindihkan.
\begin{center}
  \psset{xunit=12pt,yunit=12pt,linewidth=.4pt,arrowsize=.5pt 2.5,arrowinset=.1,
         arrowlength=1.75}
  \pspicture(-3,-2)(3,2)
    \psaxes[labels=none,linecolor=weakgray,tickstyle=bottom,
        ticks=all,ticksize=1.5pt]{<->}(0,0)(-2.75,-1.75)(2.75,1.75)
     \parametricplot[plotstyle=curve,linecolor=preimagegray,linewidth=.6pt]%
       {0}{180}{t cos t sin}
        \psline[linecolor=preimagegray]{->}(-1,0)(0,-1)
        \psline[linecolor=preimagegray]{->}(-.707,.707)(.707,-.707)
        \psline[linecolor=preimagegray]{->}(0,1)(1,0)
%        \psline[linecolor=preimagegray]{->}(.707,.707)(.707,.707)
        \psline[linecolor=preimagegray]{->}(1,0)(0,1)
     \parametricplot[plotstyle=curve,linewidth=.6pt]%
       {0}{180}{t sin t cos}
  \endpspicture
  \qquad
  \raisebox{16pt}{$\colvec{x \\ y}\mapsto\colvec{y \\ x}$}
\end{center}
Berikut adalah geseran yang ditampilkan sebelumnya.
\begin{center}
  \psset{xunit=12pt,yunit=12pt,linewidth=.4pt,arrowsize=.5pt 2.5,arrowinset=.1,
         arrowlength=1.75}
  \pspicture(-3,-2)(3,3)
    \psaxes[labels=none,linecolor=axisgray,tickstyle=bottom,
        ticks=all,ticksize=1.5pt]{<->}(0,0)(-2.75,-1.75)(2.75,2.75)
     \parametricplot[plotstyle=curve,linecolor=preimagegray,linewidth=.6pt]%
       {0}{180}{t cos t sin}
        \psline[linecolor=preimagegray]{->}(-1,0)(-1,-2)
        \psline[linecolor=preimagegray]{->}(-.707,.707)(-.707,-.707)
%        \psline[linecolor=preimagegray]{->}(0,1)(0,1)
        \psline[linecolor=preimagegray]{->}(.707,.707)(.707,2.121)
        \psline[linecolor=preimagegray]{->}(1,0)(1,2)
     \parametricplot[plotstyle=curve,linewidth=.6pt]%
       {0}{180}{t cos 2 t cos mul t sin add}
  \endpspicture
  \qquad
  \raisebox{14pt}{$\colvec{x \\ y}\mapsto\colvec{x \\ 2x+y}$}
\end{center}
Bandingkan gambar pengaruh pemetaan ini pada persegi satuan dengan gambar ini.

Berikut pemetaan yang agak lebih rumit (fungsi koordinat keduanya sama dengan
pemetaan pada gambar sebelumnya, tetapi fungsi koordinat pertamanya berbeda).
\begin{center}
  \psset{xunit=12pt,yunit=12pt,linewidth=.4pt,arrowsize=.5pt 2.5,arrowinset=.1,
         arrowlength=1.75}
  \pspicture(-3,-2)(3,3)
    \psaxes[labels=none,linecolor=axisgray,tickstyle=bottom,
        ticks=all,ticksize=1.5pt]{<->}(0,0)(-2.75,-1.75)(2.75,2.75)
     \parametricplot[plotstyle=curve,linecolor=preimagegray,linewidth=.6pt]%
       {0}{180}{t cos t sin}
        \psline[linecolor=preimagegray]{->}(-1,0)(-1,-2)
        \psline[linecolor=preimagegray]{->}(-.707,.707)(0,-.707)
        \psline[linecolor=preimagegray]{->}(0,1)(1,1)
        \psline[linecolor=preimagegray]{->}(.707,.707)(1.414,2.121)
        \psline[linecolor=preimagegray]{->}(1,0)(1,2)
     \parametricplot[plotstyle=curve,linewidth=.6pt]%
       {0}{180}{t cos t sin add 2 t cos mul t sin add}
  \endpspicture
  \qquad
  \raisebox{14pt}{$\colvec{x \\ y}\mapsto\colvec{x+y \\ 2x+y}$}
\end{center}
Perhatikan bahwa beberapa vektor mengalami dilatasi sekaligus rotasi sebesar
suatu sudut,
\begin{equation*}
  \colvec{x \\ 2x}\mapsto\colvec{3x \\ 4x},
\end{equation*}
sedangkan vektor-vektor lain hanya mengalami dilatasi tanpa rotasi sama sekali.
\begin{equation*}
  \colvec{x \\ \sqrt{2}x}
  \mapsto
  (1+\sqrt{2})\colvec{x \\ \sqrt{2}x}
\end{equation*}

\begin{exercises}
  \item Tunjukkan pengaruh setiap matriks pada setengah lingkaran satuan bagian
    atas.
    \begin{exparts*}
       \partsitem
         $\begin{pmatrix}
            1  &1    \\
            2  &2
          \end{pmatrix}$
       \partsitem
         $\begin{pmatrix}
            2  &3    \\
            1  &1
          \end{pmatrix}$
       \partsitem
         $\begin{pmatrix}
            2  &3    \\
            1  &-1
          \end{pmatrix}$
     \end{exparts*}
     Vektor-vektor manakah yang tetap berada pada garis yang sama melalui titik
     asal?
\end{exercises}
\index{geometri vektor eigen|)}
\end{input}
